\section{Equivalences}

\bdefn

    Let $F,G\colon\C\to\D$ be functors and $\eta\colon F\To G$ be a natural transformation.
    \benum
        \item Let $H\colon\D\to\E$ be another functor, then $H\eta\colon HF\To HG$ defined by $(H\eta)_c=H(\eta_c)$ is a natural transformation.
        \item Let $K\colon\E\to\C$ be another functor, then $\eta H\colon FH\to GH$ defined by $(\eta H)_e=\eta_{Hc}$ is a natural transformation.
    \eenum
    These operations on natural transformations are called {\emphcolor whiskering}.

\edefn

\bdefn

    Let $\eta\colon F\To G$ and $\epsilon\colon G\To H$ be natural transformations between functors $\C\to\D$, then their {\emphcolor vertical composition} is a natural
    transformation $\epsilon\circ\eta\colon F\To H$ defined to be $(\epsilon\circ\eta)_X=\epsilon_X\circ\eta_X$.
    This is called {\it vertical\/} composition as we can view $F,G,H$ as being stacked vertically and $\eta,\epsilon$ as arrows pointing vertically between them, their vertical
    composition is taken by following these arrows.

    Similarly if $\eta\colon F\To G$ is natural between functor $\C\to\D$ and $\epsilon\colon J\To K$ between functors $\D\to\E$, then their {\emphcolor horizontal composition}
    is a natural transformation $\epsilon*\eta\colon JF\To KG$ defined by $(\epsilon*\eta)_X=\epsilon G\circ J\eta=K\eta\circ\epsilon F$.

\edefn

Note that for natural transformations $\alpha,\alpha',\beta,\beta'$ we have the {\bf interchange law}:
$$ (\beta'\circ\alpha')*(\beta\circ\alpha) = (\beta'*\beta)\circ(\alpha'*\alpha) $$

\bdefn

    An {\emphcolor equivalence} of two categories $\C,\D$ is a pair of functors $F\colon\C\tof\D\colon G$ such that the compositions are naturally isomorphic to the identity functors on the
    respective categories.
    That is, $GF\cong1_\C$ and $FG\cong1_\D$.
    Two categories are {\emphcolor equivalent} if there exists an equivalence between them.

\edefn

Clearly the equivalence of categories forms an equivalence relation; transitivity is due to the composition of equivalences being an equivalence.

\bdefn

    A functor $F\colon\C\to\D$ is:
    \benum
        \item {\emphcolor full} if for every $x,y\in\C$, the map $\C(x,y)\to\D(Fx,Fy),f\mapsto Ff$ is surjective;
        \item {\emphcolor faithful} if for every $x,y\in\C$, the map $\C(x,y)\to\D(Fx,Fy)$ is injective;
        \item {\emphcolor essentially surjective} if for every $d\in\D$ there exists a $c\in\C$ such that $Fc$ is isomorphic to $d$.
    \eenum
    A full and faithful functor is called {\emphcolor fully faithful}.

\edefn

\bthrm

    A functor in an equivalence is fully faithful and essentially surjective.
    Conversely any fully faithful and essentially surjective function forms an equivalence.

\ethrm

\bproof

    $F$ being fully faithful means that $\C(x,y)\to\D(Fx,Fy)$ is a bijection.
    Indeed, if $G$ is the other component in the equivalence then we have a composition of maps
    $$ \C(x,y)\xto{\ F\ }\D(Fx,Fy)\xto{\ G\ }\C(GFx,GFy)\cong\C(x,y) $$
    where the last bijection is induced by the natural isomorphism $\epsilon\colon GF\cong1_\C$.
    This maps $f\colon x\to y$ to $\epsilon_YGFf\epsilon^{-1}_X$, which is just equal to $f$ by naturality.
    Composition in the other direction gives the identity as well, so our map is bijective as desired.
    Essential surjectivity is clear: $FGx\cong x$ for all $x\in\D$.

    For the converse we construct $G\colon\D\to\C$ as the other component in the equivalence.
    Let $d\in\D$, then by essential surjectivity there exists $x\in\C$ such that $Fx\cong d$.
    Let us choose such an $x$, and define $Gd=x$.
    And for $g\colon d\to d'$, as we have $\alpha\colon d\cong Fx,\alpha'\colon d'\cong Fx'$ then $\alpha'g\alpha^{-1}\colon Fx\to Fx'$.
    By full faithfulness, there is a unique $f\colon x\to x'$ such that $Ff=g$, so define $Gg=f$.
    This can be seen to be a functor, and we can see that $F,G$ forms an equivalence.
    \qed

\eproof

\bprop

    Fully faithful functors reflect isomorphisms: if $F\colon\C\to\D$ is fully faithful and $Ff$ is an isomorphism, then so is $f$.

\eprop

\bproof

    Suppose $f\colon c\to d$ such that $Ff$ is an isomorphism.
    Then there is a unique $g\colon d\to c$ such that $Fg$ is its inverse, by full faithfulness.
    Then we have that $F(fg),F(gf)$ are both the identities, but since $F$ is faithful this means that $fg$ and $gf$ must be the identities, so $f^{-1}=g$.
    \qed

\eproof

\bdefn

    A category is {\emphcolor connected} if every finite full subcategory is connected as an undirected graph.
    That is to say, for every two objects $X,Y$ there is a finite sequence of objects $X=X_0,\dots,X_n=Y$ such that $\hom(X_i,X_{i+1})$ or $\hom(X_{i+1},X_i)$ is non-empty.

\edefn

\bcoro

    Any connected groupoid is equivalent to the automorphism group of any of its objects.

\ecoro

\bproof

    A connected groupoid is one for which there exists isomorphisms between any two objects.
    Let $g\in\c G$ and let $G=\c G(g,g)$ denote its automorphism group.
    Then the functor $\B G\to\c G$ mapping the unique object of $\B G$ to $g$ and the elements of $G$ to their respective elements in $G=\c G(g,g)$ is fully faithful and essentially surjective
    (as all objects in $\c G$ are isomorphic).
    Thus this forms an equivalence.
    \qed

\eproof

\bdefn

    Let $\C$ be a category, its {\emphcolor skeleton} is the subcategory $\sk\C\subseteq\C$ obtained by choosing a single object from each isomorphism class in $\C$.

\edefn

Note that we say {\it the\/} skeleton of $\C$ even though the construction of $\sk\C$ depends on choosing a representative from each isomorphism class.
But it is not hard to see that any choice gives isomorphic categories, thus $\sk\C$ is defined uniquely up to isomorphism.

\bprop

    The inclusion $\sk\C\to\C$ is an equivalence of categories.

\eprop

\bproof

    This is clear, since $\sk\C$ is a full subcategory the inclusion is fully faithful.
    By definition of $\sk\C$ the inclusion is also essentially surjective, as desired.
    \qed

\eproof

For a connected groupoid, we clearly have that $\sk\c G$ is simply a category with a single object whose hom-set is the automorphism group of any object in $\c G$.
That is, $\sk\c G=\B\c G(g,g)$ for any $g\in\c G$.
So our previous claim is just a special example of the above proposition.

\bexam

    Let $\C$ be a category equivalent to the trivial category $[0]$.
    Then $F\colon\C\to[0]$ is fully faithful, which means that $\hom(x,y)\cong\hom_0(0,0)$, i.e.\ it is a singleton.
    So between every two objects in $\C$ there is a unique morphism, in particular it must be an isomorphism.

    In general if $\C$ is equivalent to a groupoid, it is a groupoid.
    Notice that if $F\colon\C\to\D$ is an equivalence then every $\D$-morphism is of the form $\alpha\circ Ff\circ\beta$ for $\alpha,\beta$ isomorphisms and $f$ a $\C$-morphism.
    So if $F\colon G\to\C$ is an equivalence then every $\C$-morphism is of the form $\alpha\circ Fg\circ\beta$, and since $g$ is an isomorphism so is $Fg$, so every $\C$-morphism is an isomorphism.

\eexam

\bprop

    For categories $\C,\D$ the following are equivalent:
    \benum
        \item $\C$ and $\D$ are equivalent,
        \item $\sk\C$ and $\sk\D$ are isomorphic,
        \item $\sk\C$ and $\sk\D$ are equivalent.
    \eenum

\eprop

\bproof

    If $F\colon\sk\C\to\sk\D$ is an equivalence, then it must be surjective by essential surjectivity.
    Furthermore it must be injective: if $F(c)=F(c')$ then in particular $F(c),F(c')$ are isomorphic and since $F$ reflects isomorphisms, $c\cong c'$.
    Thus $c=c'$ since the skeleton has a unique object of each isomorphism class.
    This means that $F$ is an isomorphism, so $(2)$ and $(3)$ are equivalent.

    If $F\colon\C\to\D$ is an equivalence then it restricts to an equivalence $\sk\C\to\sk\D$ (as it is still fully faithful and essentially surjective).
    And if $\sk\C$ and $\sk\D$ are equivalent, since $\C$ and $\D$ are equivalent to their skeletons, by transitivity $\C$ and $\D$ are equivalent.
    \qed

\eproof

Thus if $\C$ is equivalent to $[0]$, then $\sk\C\cong\sk[0]=[0]$, i.e.\ it has a single isomorphism class, as we already showed.
And if $\C$ is equivalent to a groupoid $\c G$, then $\sk\C\cong\sk\c G$ which has a single object whose hom-set is only isomorphisms.
Thus $\C$ must be a groupoid: it has a single isomorphism class and all of its morphisms are isomorphisms.

\bexam

    \def\Vec{{\rm Vec}}\def\fin{{\sf fin}}\def\mat{{\rm Mat}}
    Let $k$ be a field, then we have a category $\Vec_k^\fin$ of finite-dimensional $k$-vector spaces.
    We also have $\mat_k$, the category of matrices over $k$ defined as follows:
    \benum
        \item Objects in $\mat_k$ are natural numbers (including $0$),
        \item Morphisms $n\to m$ are matrices of dimension $m\times n$, i.e.\ $\mat_k(n,m)=\mat_{m\times n}(k)$,
        \item For matrices $A\colon n\to m$ and $B\colon m\to p$, define their composition to be their multiplication: $B\circ A=BA$.
    \eenum

    These categories are equivalent: we define a functor $F\colon\mat_k\to\Vec_k^\fin$ by $F(n)=k^n$ and for $A\colon n\to m$, $F(A)\colon k^n\to k^m,v\mapsto Av$.
    Simple linear algebra shows that this is essentially surjective and fully faithful.

    Note that $\sk\mat_k=\mat_k$, and thus by equivalence $\sk\Vec_k^\fin\cong\mat_k$.

\eexam

\bexam

    We define the category of sets with partial functions $\Set_p$ as follows:
    \benum
        \item Objects are sets,
        \item Morphisms $X\to Y$ are partial functions from $X$ to $Y$, which we view as pairs $(U,f)$ where $U\subseteq X$ and $f\colon U\to Y$,
        \item Composition is defined as follows: for $(U,f)\colon X\to Y$ and $(V,g)\colon Y\to Z$, $(V,g)\circ(U,f)=(f^{-1}(V),g\circ f)$.
    \eenum

    This is equivalent to $\Set_*$, the category of pointed sets (objects are $(X,x)$ for $x\in X$ and morphisms $f\colon(X,x)\to(Y,y)$ are functions $X\to Y$ such that $f(x)=y$).
    Indeed, we define a functor $F\colon\Set_p\to\Set_*$ by $F(X)=(X\dcup\Set*,*)$ and for $(U,f)\colon X\to Y$ define $F(U,f)\colon X\dcup\Set*\to Y\dcup\Set*$ by $*\mapsto*$, $x\mapsto f(x)$ if defined and
    otherwise to $*$.
    This is essentially surjective clearly, and fully faithful clearly.

\eexam

\bdefn

    For a category $\C$, we define its {\emphcolor center} to be the set of natural transformations $1_\C\To1_\C$.
    Denote this by $Z(\C)$.

\edefn

Note that $Z(\C)$ equipped with composition (vertical and horizontal agree on the center) is a commutative monoid.
Clearly it is a monoid, as the natural transformation defined by the identity on all components is its identity.
And it is commutative; by naturality for every $f\colon X\to X$ one has $f\alpha_X=\alpha_Xf$, in particular for $f=\beta_X$.

\bprop

    If $\C$ and $\D$ are equivalent categories, then their centers are isomorphic commutative monoids.

\eprop

\bproof

    Let $F\colon\C\tof\D\colon G$ be an equivalence, with isomorphisms $\epsilon\colon FG\cong1_D$ and $\eta\colon1_C\cong GF$.
    Then we define $Z(\C)\to Z(\D)$ as follows: $\alpha\mapsto\epsilon\circ F\alpha G\circ\epsilon^{-1}$.
    This is a monoid morphism: $(\epsilon\circ F1G\circ\epsilon^{-1})_X=\epsilon_X\circ 1_{FGx}\circ\epsilon^{-1}_X=1_X$ and since $F\alpha\beta G=(F\alpha G)(F\beta G)$ we get commutativity.
    Its inverse is given by $\beta\mapsto\eta^{-1}\circ G\beta F\circ\eta$.
    \qed

\eproof

